\chapter{Differentiable structure (planned)}
\label{ch:c8}

\begin{keybox}\textbf{PLANNED --- not yet written.} This chapter and its
module (\file{08\_differentiable\_structure/}) are a placeholder. The
outline below is what a future session will build from; see
\file{HANDOFF.md} for the full instructions.\end{keybox}

\emph{Module: \file{08\_differentiable\_structure/} (stub). The NS-SBM
engine as an instrument for design: gradients through the solve, and shape
optimization of the immersed body.}

\section*{What this chapter will teach}
\begin{itemize}
\item \textbf{The adjoint of the NS-SBM engine.} A gradient of a flow QoI
  (drag, a wake observable) with respect to a control --- the SBM shift /
  the body shape, the Reynolds number, a boundary trace. The steady adjoint
  is a \emph{single transposed solve} (no trajectory), which is exactly what
  makes it tractable at 100M DOF.
\item \textbf{Three-way-verified gradients.} Following the sibling
  differentiable track's discipline: custom adjoint $==$ autograd twin $==$
  finite difference, on the \emph{same} operator the forward stepper
  marches. The course's standing bar is that a gradient is not trusted
  until all three agree.
\item \textbf{Shape optimization on the shifted boundary.} Because the SBM
  shift is fixture \emph{data} ($d$, $\text{corr}$; Ch.~\ref{ch:c4}),
  differentiating with respect to the body geometry flows through the
  distance field --- a natural, mesh-free shape-design surface.
\end{itemize}

\section*{Source material (for the author of this chapter)}
\begin{itemize}
\item The SBM adjoint / shape bricks already in the tree:
  \file{src/diffsim/sbm/ns\_adjoint.py}, \file{leray\_adjoint.py},
  \file{ns\_shape.py}, \file{transient\_adjoint.py}, \file{adjoint.py}.
\item The SP-1 R1 XDD adjoint work (MERGED \code{fbfeed5}): the
  three-way-verification pattern (adjoint vs torch-twin vs FD), the unified
  \code{Control} interface, the exposed sensitivity rows --- the template
  this chapter mirrors for the NS-SBM operator.
\item The P2 milestone's ``steady adjoint = 1 transposed solve, tractable at
  100M'' argument (deliverable C, the maize steady-differentiable design
  engine).
\item The sibling \file{orgelmorph-course/differentiable/} track for the
  module template of a gradient chapter (core module, \code{run.py} printing
  the three-way agreement, \code{EXPECTED.md}).
\end{itemize}

\section*{Status}
Future. Ties into the E-differentiable track and the SP-1 R1 adjoint work;
best built after the forward engine chapters (00--05) and the scaling
chapter (06) are the students' shared baseline.
